% Options for packages loaded elsewhere
\PassOptionsToPackage{unicode}{hyperref}
\PassOptionsToPackage{hyphens}{url}
\documentclass[
  11pt,
  a4paper,
]{article}
\usepackage{xcolor}
\usepackage[margin=1in]{geometry}
\usepackage{amsmath,amssymb}
\setcounter{secnumdepth}{-\maxdimen} % remove section numbering
\usepackage{iftex}
\ifPDFTeX
  \usepackage[T1]{fontenc}
  \usepackage[utf8]{inputenc}
  \usepackage{textcomp} % provide euro and other symbols
\else % if luatex or xetex
  \usepackage{unicode-math} % this also loads fontspec
  \defaultfontfeatures{Scale=MatchLowercase}
  \defaultfontfeatures[\rmfamily]{Ligatures=TeX,Scale=1}
\fi
\usepackage{lmodern}
\ifPDFTeX\else
  % xetex/luatex font selection
  \setmainfont[]{Latin Modern Roman}
\fi
% Use upquote if available, for straight quotes in verbatim environments
\IfFileExists{upquote.sty}{\usepackage{upquote}}{}
\IfFileExists{microtype.sty}{% use microtype if available
  \usepackage[]{microtype}
  \UseMicrotypeSet[protrusion]{basicmath} % disable protrusion for tt fonts
}{}
\makeatletter
\@ifundefined{KOMAClassName}{% if non-KOMA class
  \IfFileExists{parskip.sty}{%
    \usepackage{parskip}
  }{% else
    \setlength{\parindent}{0pt}
    \setlength{\parskip}{6pt plus 2pt minus 1pt}}
}{% if KOMA class
  \KOMAoptions{parskip=half}}
\makeatother
\usepackage{longtable,booktabs,array}
\newcounter{none} % for unnumbered tables
\usepackage{calc} % for calculating minipage widths
% Correct order of tables after \paragraph or \subparagraph
\usepackage{etoolbox}
\makeatletter
\patchcmd\longtable{\par}{\if@noskipsec\mbox{}\fi\par}{}{}
\makeatother
% Allow footnotes in longtable head/foot
\IfFileExists{footnotehyper.sty}{\usepackage{footnotehyper}}{\usepackage{footnote}}
\makesavenoteenv{longtable}
\setlength{\emergencystretch}{3em} % prevent overfull lines
\providecommand{\tightlist}{%
  \setlength{\itemsep}{0pt}\setlength{\parskip}{0pt}}
\usepackage{mathtools}
\usepackage{enumitem}
\usepackage{microtype}
\usepackage{booktabs}
\usepackage{array}
\usepackage{bookmark}
\IfFileExists{xurl.sty}{\usepackage{xurl}}{} % add URL line breaks if available
\urlstyle{same}
\hypersetup{
  pdftitle={Quantum Computation in Event Density A Substrate-Level Architecture and Its Limits},
  pdfauthor={Allen Proxmire},
  hidelinks,
  pdfcreator={LaTeX via pandoc}}

\title{Quantum Computation in Event Density\\
\large A Substrate-Level Architecture and Its Limits}
\author{Allen Proxmire}
\date{May 2026}

\begin{document}
\maketitle

\subsection{Abstract}\label{abstract}

The Event Density (ED) framework provides a substrate-level account of
quantum computation in which three substrate quantities --- local
\textbf{multiplicity} (the count of available participation pathways at
the system's endpoints), \textbf{cross-endpoint connectivity} (the
substrate-mediated bandwidth that links them), and
\textbf{commitment-injection rate} (the rate at which the environment
forces individuation) --- jointly determine whether quantum computation
is possible. The \textbf{Unresolved-Regime Characterization Theorem
(UR-1)} identifies three independently-necessary conditions on these
quantities, and the resulting QC operating window is given structurally
by
\(\tau_{\mathrm{QC}} = \min(\tau_{(\mathrm{i})}, \tau_{(\mathrm{ii})}, \tau_{(\mathrm{iii})})\)
over the three corresponding failure timescales. Each architectural
platform binds on whichever timescale is shortest under its protection
strategy.

Three architectural classes --- \textbf{engineered-low-multiplicity (A)}
covering superconducting qubits, trapped ions, and gate-model photonics;
\textbf{global-geometric-rigidity (B)} covering topological qubits and
photonic Chern channels; and \textbf{high-multiplicity-redundancy (C)}
covering bosonic codes and multi-timescale photonics --- exhaust the
substrate-allowed strategies for holding the operational regime. Error
correction, dynamical decoupling, reservoir engineering, and hybrid
architectures are derived as compositions or extensions of the three
classes rather than new strategies. The \textbf{multiplicity-cap
function \(M\)} unifies the classes as three projections of one
substrate object, with platform-specific behavior obtained by
specializing class-dependent parameters.

The framework's strongest output is \textbf{cross-platform unification}:
the matter-wave quantum-classical boundary observed empirically at
140--250 kDa molecular mass and the multiplicity walls that limit
qubit-system scaling are the same substrate-determined boundary
projected onto two different platforms, evaluated at consistent values
of one underlying substrate constant. Sharp predictions follow directly:
superconducting platforms encounter a hard scaling ceiling unless they
are composed with logical-qubit encoding; topological qubits can outrun
this ceiling exponentially in the topological gap if topology
engineering can deliver stable structures; redundancy-based
architectures saturate at a substrate-determined correlation budget.
Four classes of near-term experimental tests with explicit falsification
conditions are specified.

The structural form of the framework is derived from substrate
primitives; specific numerical thresholds (the critical multiplicity,
the correlation budget, the maximum stable topological gap, and the
platform-specific wall locations) are downstream empirical calibrations
rather than free parameters, with the empirical anchors that fix them
identified explicitly. \textbf{The verdict is architectural-level
structural completeness with an explicit empirical-calibration roadmap}:
the substrate-level account of quantum computation is delivered, the
predictions are sharp enough to be falsified by near-term measurements,
and the calibration program that converts structural predictions into
platform-specific numerical bounds is engineering-tractable rather than
theoretically blocked.

\begin{center}\rule{0.5\linewidth}{0.5pt}\end{center}

\subsection{1. Introduction}\label{introduction}

\subsubsection{1.1 Motivation}\label{motivation}

Quantum computing has produced, over the past three decades, two
compelling but partial pictures of itself. The first picture is
engineering-led: quantum computers are devices that maintain coherence
against decoherence, and progress is the steady reduction of the
engineering distance between current platforms and a fault-tolerant
million-qubit logical machine. The second picture is
information-theoretic: quantum computation is the manipulation of
entangled state-spaces whose dimension grows exponentially in qubit
count, and the question is what classes of problems benefit from the
exponential resource.

Neither picture answers the structural question: \emph{what is quantum
computation, at the level of physics, and what are its substrate-level
limits?} Empirical scaling debates --- when will logical-qubit lifetimes
exceed physical T₁? when will topological qubits surpass
superconducting? does multi-axis redundancy keep paying as you add axes?
--- are addressed pragmatically, platform by platform, with no shared
structural framework predicting which architectures should hit which
walls and why.

The Event Density program approaches this question from a
substrate-ontological starting point. ED's primitives --- local
participation density, irreversible commitment events, finite-bandwidth
substrate channels, V1 finite-width vacuum kernel --- generate quantum
mechanics as a coarse-grained continuum reading of substrate dynamics in
the thin-participation regime (the program's Phase-1 closure of the QM
postulates). Quantum computation is then a particular \emph{engineered
occupation} of this regime: a system held in the unresolved
low-multiplicity state long enough for its substrate-level participation
geometry to be manipulated, before environmental ED-injection forces
individuation at readout.

This paper develops the substrate-level architecture of that picture. It
establishes the gate-condition theorem for the unresolved regime,
derives the three-class exhaustive taxonomy of architectural protection
strategies, constructs the multiplicity-cap function \(M\) that unifies
QC platform behavior across substrate-distant scales, and delivers
falsifiable predictive content with explicit demarcation between
substrate-derived form and substrate-inherited values.

\subsubsection{1.2 ED Ontology in Brief}\label{ed-ontology-in-brief}

ED treats reality, at its most fundamental level, as discrete
\emph{events} --- atomic moments of becoming --- linked by relational
\emph{participation channels}. Particles and fields are not fundamental;
they are stable patterns of substrate participation. Spacetime is not
fundamental; it emerges from the way events relate, with locality,
causality, and metric structure as coarse-grained signatures of dense,
redundant, irreversibly-committed participation networks. Time is not a
primitive; time is the irreversible direction along which events commit
to being something, with the asymmetry secured by the substrate
primitive P11 (commitment-irreversibility).

For quantum computation, three substrate primitives are load-bearing.

\begin{itemize}
\item
  \textbf{Multiplicity} \(\mathcal{M}\): the count of viable distinct
  ED-gradient pathways available to a system. Per ED-I-01
  (Superconductivity, February 2026), multiplicity is the ED analogue of
  entropy. Low \(\mathcal{M}\) means few alternative pathways; high
  \(\mathcal{M}\) means many. Symmetry collapses multiplicity; thermal
  agitation raises it.
\item
  \textbf{Cross-bandwidth} \(\Gamma_{\mathrm{cross}}\): the
  substrate-mediated rate at which adjacent regions exchange correlated
  participation events. Established in Arc D (Diffusion Coarse-Graining
  Theorem, April 2026) as
  \(\Gamma_{\mathrm{cross}} \sim \exp[-\alpha\sigma]\) where
  \(\sigma = |\nabla\rho|\ell_P^2/\rho_{\mathrm{local}}\) is the
  substrate-scale gradient sparsity.
\item
  \textbf{Commitment-injection rate} \(\Lambda\): the rate at which P11
  commitment events fire at a region's endpoints, sourced by
  environmental ED-injection plus internal participation dynamics.
\end{itemize}

A system performing quantum computation, in ED's reading, is a substrate
region in which a participation rule remains \emph{unresolved} across
designated endpoints --- its identity has not individuated to any
specific configuration --- for a finite substrate-time window long
enough for the rule's geometry to be manipulated through engineered
substrate operations (gates), before environmental ED-injection drives
it to individuation (readout).

\subsubsection{1.3 Why Standard QC Framing Misses the Substrate-Level
Question}\label{why-standard-qc-framing-misses-the-substrate-level-question}

Standard QC treats coherence as a state property of a wavefunction and
decoherence as the wavefunction's coupling to environmental degrees of
freedom. The ontological picture is: there is a quantum system, it has a
state, the state can be coherent or decoherent, and engineering reduces
the coupling that drives state collapse.

ED's reframe is structurally distinct. There is no fundamental
wavefunction in the substrate ontology; what exists is a participation
structure with a rule spanning endpoints. Coherence is the integrity of
that rule before commitment; decoherence is the substrate-level event of
commitment occurring. Engineering doesn't slow down a coupling-rate to a
state; it holds the substrate region in a low-multiplicity,
well-connected, low-commitment-rate configuration, against the
substrate's natural tendency to proliferate gradients and force
individuation.

Three structural questions become tractable in this picture that are not
addressable in the standard one:

\begin{enumerate}
\def\labelenumi{\arabic{enumi}.}
\item
  \emph{What is the substrate-level signature of ``doing QC''?} Answer
  (this paper): a system satisfying the three conditions of the
  Unresolved-Regime Characterization Theorem (UR-1).
\item
  \emph{What architectural strategies can hold the system in that
  regime?} Answer: three exhaustive substrate-level protection
  strategies --- direct \(\mathcal{M}\) control (Class A), direct
  \(\sigma\)-geometry control via topology (Class B), redundancy across
  pathways (Class C).
\item
  \emph{What unifies coherence behavior across qualitatively different
  platforms --- matter-wave interferometers, superconducting qubits,
  Josephson junctions, topological qubits, multi-timescale photonic
  lattices?} Answer: one substrate function \(M\) with three
  architectural-class projections, on which all platforms place at
  consistent substrate-determined locations.
\end{enumerate}

\subsubsection{1.4 The Conceptual Shift: From Decoherence to
Multiplicity}\label{the-conceptual-shift-from-decoherence-to-multiplicity}

The reframing this paper develops is not a refinement of the standard
quantum-computing picture; it is a substitution at the level of which
substrate quantity is the load-bearing variable.

In the standard picture, \emph{decoherence} is the central object --- a
coupling-rate between a quantum state and its environment, to be reduced
by isolation engineering, dynamical decoupling, error-correcting codes,
and similar techniques. Coherence times are explained by where the
coupling-rate sits; scaling limits are explained by where the
coupling-rate becomes irreducible.

In the substrate picture developed here, \emph{multiplicity} is the
central object --- a substrate-state quantity that counts available
ED-gradient pathways, with low-multiplicity regions admitting the
unresolved-rule regime that quantum computation requires and
high-multiplicity regions forcing individuation. Decoherence in the
standard sense is the \emph{coarse-grained signature} of substrate-level
individuation events firing at endpoints whose multiplicity has risen
above a substrate-determined threshold. The standard variable is real
but downstream; the substrate variable is upstream and structurally
prior.

This framework replaces the standard decoherence-centric view with a
substrate-level multiplicity-centric account, unifying matter-wave
interference, superconducting qubits, Josephson-junction devices,
topological qubits, multi-timescale photonic platforms, and
redundancy-based bosonic codes under a single substrate mechanism ---
the multiplicity-cap function \(M\). The reframing has direct
implications for how quantum-computing scaling limits are interpreted
and pursued: limits that appear in the standard picture as engineering
accidents (residual environmental couplings, fabrication imperfections,
correlated noise channels) appear in the substrate picture as
substrate-determined structural ceilings whose locations are constrained
by the cross-platform identity of \(M\).

The paper does not claim that engineering effort is misdirected ---
engineering against the substrate-level binding constraints is exactly
the right work in the appropriate regime. It claims that the
\emph{interpretive frame} under which engineering progress is evaluated
is shifted: progress against substrate-level walls is structurally
bounded; progress that addresses non-binding constraints is wasted; and
the cross-platform identity of substrate-level walls means that progress
in one domain (e.g., matter-wave-interferometry mass-cap measurement)
directly constrains what is achievable in another (e.g.,
superconducting-qubit system-multiplicity scaling).

\subsubsection{1.5 Structure of the Paper}\label{structure-of-the-paper}

Section 2 develops the substrate primitives and DCGT machinery. Section
3 derives UR-1. Section 4 quantifies failure-mode rates. Section 5
audits the three architectural classes for substrate-level
exhaustiveness. Section 6 constructs the multiplicity-cap function
\(M\). Section 7 delivers cross-platform unification and predictive
content. Section 8 provides FORCED / INHERITED / OPEN classification.
Section 9 is discussion and implications. Appendices and metadata
follow.

\begin{center}\rule{0.5\linewidth}{0.5pt}\end{center}

\subsection{2. Substrate Primitives and DCGT
Machinery}\label{substrate-primitives-and-dcgt-machinery}

\subsubsection{2.1 Participation Density and
Gradients}\label{participation-density-and-gradients}

The substrate's load-bearing scalar field is the local participation
density \(\rho(\mathbf{x})\), defined as the density of viable
participation-channel slots per coarse-graining cell of size
\(R_{\mathrm{cg}}\). Gradients \(\nabla\rho\) encode spatial variation.
The dimensionless ratio \[
\sigma(\mathbf{x}) \equiv \frac{|\nabla\rho|\,\ell_P^2}{\rho_{\mathrm{local}}}
\] measures the substrate-scale steepness of \(\rho\) at the Planck
length \(\ell_P\). This is the load-bearing substrate quantity for
cross-bandwidth derivation in Arc D and for the present arc's condition
(ii) analysis.

\subsubsection{2.2 V1 Finite-Width Kernel}\label{v1-finite-width-kernel}

The V1 vacuum kernel, established in Theorem N1 (April 2026) and refined
in Theorem 18 (April 2026), is the finite-width temporal smearing kernel
that mediates substrate participation events. Its width sets the
temporal resolution of commitment events at a region. V1 carries the
kernel-level arrow of time (Theorem 18): the kernel is
forward-cone-only, propagating substrate correlations forward in
P11-time only.

\subsubsection{2.3 Multi-Scale Expansion and the Hydrodynamic
Window}\label{multi-scale-expansion-and-the-hydrodynamic-window}

DCGT (Diffusion Coarse-Graining Theorem, Arc D, April 2026) establishes
a hydrodynamic-window scale separation \[
\ell_P \ll R_{\mathrm{cg}} \ll L_{\mathrm{flow}}
\] under which substrate dynamics admit a multi-scale expansion to
coarse-grained continuum equations. DCGT supplies the cross-bandwidth
structure used throughout this paper.

\subsubsection{2.4 Cross-Bandwidth and
Sparsity}\label{cross-bandwidth-and-sparsity}

For two adjacent substrate regions separated by a pathway of length
\(L\) along which \(\sigma\) varies, the cross-bandwidth is \[
\Gamma_{\mathrm{cross}}(\mathbf{x}_1, \mathbf{x}_2) \sim \exp\!\left[-\alpha\!\int_{\mathrm{path}}\sigma(\mathbf{x})\,d\ell\right]
\] with \(\alpha\) a substrate-determined dimensionless prefactor. The
exponential structure is forced by the multi-scale expansion: large
\(|\nabla\rho|\) at the substrate scale creates steep barriers to
participation-channel mediation between adjacent regions. This is the
same machinery that produces the BH-2 horizon-formation condition (Arc
BH, 2026-05-01); here it is applied at engineered-system scale rather
than gravitational-collapse scale.

\subsubsection{2.5 Multiplicity and the ED-Entropy
Analogue}\label{multiplicity-and-the-ed-entropy-analogue}

Per ED-I-01, multiplicity is the ED-substrate analogue of entropy. A
region with high \(\mathcal{M}\) supports many distinct participation
pathways; ED-flow can branch, scatter, and decohere across them. A
region with low \(\mathcal{M}\) supports few distinct pathways; ED-flow
is constrained. Symmetry, in ED, is not geometric ornamentation but a
\emph{reduction} in the number of distinct ED-gradient configurations
the system can occupy. Decoherence is re-entry into high-multiplicity
participation: when a system's substrate environment proliferates
gradient pathways, the system's identity individuates to one of them.

This identification is load-bearing for the architectural classification
in Section 5.

\subsubsection{2.6 Commitment-Irreversibility
(P11)}\label{commitment-irreversibility-p11}

P11 is the substrate primitive that participation events, once
committed, cannot be uncommitted. Commitment is the moment a chain
``becomes'' something it wasn't a moment before. P11 is the only
direction-bearing substrate primitive in ED; every other primitive (V1
kernel, participation density, cross-bandwidth) is time-symmetric in
isolation. The kernel-level arrow of time (Theorem 18) propagates from
P11 through the V1 kernel to produce forward-cone-only causal structure
at substrate level.

For QC: P11 fires every time the system commits a participation event
--- every gate operation, every ancilla reset, every measurement, every
internal scattering event. The cumulative commitment density inside a
system bounds how long the unresolved-rule regime can be held.

\begin{center}\rule{0.5\linewidth}{0.5pt}\end{center}

\subsection{3. UR-1: The Unresolved-Regime Characterization
Theorem}\label{ur-1-the-unresolved-regime-characterization-theorem}

\subsubsection{\texorpdfstring{3.1 The Functional Form of
\(\mathcal{U}\)}{3.1 The Functional Form of \textbackslash mathcal\{U\}}}\label{the-functional-form-of-mathcalu}

Define \(\mathcal{U}(\mathcal{S},t) \in [0,1]\) as the integrity of an
unresolved participation rule across the designated endpoints of
substrate region \(\mathcal{S}\) at substrate time \(t\):
\(\mathcal{U} \to 1\) when the rule remains globally coherent across all
endpoints (no individuation has occurred); \(\mathcal{U} \to 0\) when
full individuation has been reached.

Three substrate processes can degrade \(\mathcal{U}\), each from a
structurally distinct substrate quantity. They compose multiplicatively,
because any one alone destroys \(\mathcal{U}\) regardless of the others:
\[
\mathcal{U}(\mathcal{S},t) = \prod_{i \in \mathrm{endpts}(\mathcal{R})} \mu\!\left(\frac{\mathcal{M}_i(t)}{\mathcal{M}_{\mathrm{crit}}}\right) \cdot \prod_{(i,j) \in \mathcal{R}} \kappa\!\left(\frac{\gamma_{ij}(t)}{\Gamma_{\mathrm{min}}}\right) \cdot \exp\!\left[-\int_0^t \Lambda_{S}(t')\,dt'\right].
\]

The three factors correspond, in order, to \emph{multiplicity headroom},
\emph{rule-spanning connectivity}, and \emph{commitment-survival}.

The form is \textbf{FORCED} by substrate machinery:

\begin{itemize}
\tightlist
\item
  \(\mu, \kappa : [0,\infty) \to [0,1]\) are monotone functions. \(\mu\)
  decreasing from 1 to 0 (more multiplicity = less unresolvedness
  preservation); \(\kappa\) increasing from 0 to 1 (more cross-bandwidth
  = better rule-spanning). Specific shapes (Boltzmann, rational,
  sigmoid) are \textbf{INHERITED} from V1-kernel + DCGT closed-form
  details.
\item
  The exponential survival factor is forced by the Poisson-class
  structure of P11 commitment events: if \(\Lambda_{S}(t')\) is the
  local rate of individuation-driving commitments, the probability that
  no such event has fired in \([0,t]\) is exactly
  \(\exp[-\int \Lambda dt']\).
\item
  \(\mathcal{R}\) is the set of endpoint-pairs across which the
  participation rule spans. \(\mathcal{M}_{\mathrm{crit}}\) and
  \(\Gamma_{\mathrm{min}}\) are substrate-determined critical
  thresholds, INHERITED.
\end{itemize}

\subsubsection{3.2 The UR-1 Theorem}\label{the-ur-1-theorem}

\begin{quote}
\textbf{UR-1 (Unresolved-Regime Characterization Theorem).} Let
\(\mathcal{S}\) be a substrate region with designated participation-rule
relation set \(\mathcal{R}\) over endpoints \(\{e_1, \ldots, e_n\}\).
Let \(\mathcal{U}_{\mathrm{min}} \in (0,1)\) be a target unresolvedness
threshold. Then
\(\mathcal{U}(\mathcal{S},t) \geq \mathcal{U}_{\mathrm{min}}\) holds at
substrate time \(t\) if and only if:

\textbf{(i) Multiplicity bounded.}
\(\mathcal{M}_i(t) \leq \mathcal{M}_{\mathrm{crit}}^{(\mu^{-1}(\mathcal{U}_{\mathrm{min}}^{1/n}))}\)
for every endpoint \(e_i\).

\textbf{(ii) Cross-endpoint connectivity sustained.}
\(\gamma_{ij}(t) \geq \Gamma_{\mathrm{min}} \cdot \kappa^{-1}(\mathcal{U}_{\mathrm{min}}^{1/|\mathcal{R}|})\)
for every pair \((i,j) \in \mathcal{R}\).

\textbf{(iii) Commitment-injection bounded.}
\(\int_0^t \Lambda_{S}(t')\,dt' \leq \ln(1/\mathcal{U}_{\mathrm{min}}) - C_{(\mathrm{i})} - C_{(\mathrm{ii})}\),
where \(C_{(\mathrm{i})}, C_{(\mathrm{ii})}\) are negative-log deficits
from conditions (i) and (ii).

Each condition is independently necessary. Failure of any one drives
\(\mathcal{U}\) below \(\mathcal{U}_{\mathrm{min}}\) at a
substrate-determined rate. The conditions are coupled through substrate
geometry (substrate configurations cannot independently dial all three),
but each can fail without the others.
\end{quote}

\subsubsection{3.3 Reduction to Known
Regimes}\label{reduction-to-known-regimes}

Three known regimes recover from UR-1, each with a different binding
condition:

\textbf{Bulk superconductor (ED-I-01).} Lattice symmetry collapses local
multiplicity to \(\mathcal{M} \approx 1\) across the bulk; \(\sigma\) is
smooth, \(\gamma_{ij}\) high. Below \(T_c\), condition (iii) holds with
margin: thermal ED-injection rate is below the substrate individuation
threshold. Above \(T_c\), thermal agitation overwhelms the
symmetry-locking; condition (iii) crosses threshold; \(\mathcal{U}\)
drops; resistance reappears. \textbf{\(T_c\) is identified as the
substrate temperature at which condition (iii) crosses its threshold.}

\textbf{Josephson junction (ED-I-23).} Two low-\(\mathcal{M}\) SC
electrodes separated by a thin insulating barrier. (i) holds in both
electrodes; (iii) holds with adequate isolation. Condition (ii) is the
load-bearing constraint by design: the barrier is engineered such that
\(\gamma_{LR}\) is just above \(\Gamma_{\mathrm{min}}\). \textbf{JJ MQT
(Devoret-Martinis-Clarke 1985) corresponds to dynamical fluctuations at
the (ii) boundary} --- the rule globally reconfigures across the gap
when \(\gamma_{LR}\) momentarily drops below threshold. The MQT rate's
WKB exponential structure is recovered directly from DCGT
\(\kappa\)-function evaluation: \[
\tau_{\mathrm{MQT}}^{-1} \sim \omega_0\,\exp\!\left[-\alpha\!\int_{\mathrm{barrier}}\sigma\,d\ell\right].
\]

\textbf{Free atomic-scale matter wave.} The molecule is itself a
low-\(\mathcal{M}\) ED-object whose internal \(\mathcal{M}\) scales with
mass (rotational + vibrational + electronic DOF activation). For low
molecular mass, \(\mathcal{M} \ll \mathcal{M}_{\mathrm{crit}}\) and (i)
holds. For high mass, internal \(\mathcal{M}\) rises; at the Q-C
boundary (\(M_\star \sim 140\)--\(250\) kDa empirically), internal
\(\mathcal{M}\) crosses \(\mathcal{M}_{\mathrm{crit}}\). \textbf{The
matter-wave Q-C boundary is identified as the molecular mass at which
condition (i) crosses threshold.}

Three different empirical phenomena, three different UR-1 conditions,
one substrate framework.

\begin{center}\rule{0.5\linewidth}{0.5pt}\end{center}

\subsection{4. Failure-Mode Rates and the QC Operating
Window}\label{failure-mode-rates-and-the-qc-operating-window}

Each UR-1 condition admits a substrate-level dynamical equation
governing the timescale at which it crosses threshold.

\subsubsection{4.1 Condition (i): Multiplicity
Dynamics}\label{condition-i-multiplicity-dynamics}

\[
\frac{d\mathcal{M}}{dt}\bigg|_{S} = \alpha_{\mathrm{env}}\,\Lambda_{\mathrm{env}}(t) + \alpha_{\mathrm{act}}\,S_{\mathrm{int}}(t) - A_{S}\,(\mathcal{M} - \mathcal{M}_{\mathrm{floor}})
\]

with:

\begin{itemize}
\tightlist
\item
  \(\alpha_{\mathrm{env}}\,\Lambda_{\mathrm{env}}\): environmental
  ED-injection driving pathway-count growth.
\item
  \(\alpha_{\mathrm{act}}\,S_{\mathrm{int}}\): internal
  pathway-activation rate from thermal or operational stimulation.
\item
  \(A_{S}\,(\mathcal{M} - \mathcal{M}_{\mathrm{floor}})\): architectural
  restoring rate, returning \(\mathcal{M}\) to the
  architecturally-imposed floor \(\mathcal{M}_{\mathrm{floor}}\).
\end{itemize}

Two failure modes:

\begin{itemize}
\tightlist
\item
  \textbf{Static failure:}
  \(\mathcal{M}_{\mathrm{floor}} \geq \mathcal{M}_{\mathrm{crit}}\). The
  system cannot be operated regardless of timescale.
  \(\tau_{(\mathrm{i})} = 0\). \emph{This is the matter-wave Q-C
  boundary regime.}
\item
  \textbf{Dynamic failure:}
  \(\mathcal{M}_{\mathrm{floor}} < \mathcal{M}_{\mathrm{crit}}\) but
  environmental injection drives steady-state above threshold. Finite
  \(\tau_{(\mathrm{i})}\) set by net multiplicity-growth rate after
  architectural restoring.
\end{itemize}

\subsubsection{4.2 Condition (ii): Cross-Bandwidth
Dynamics}\label{condition-ii-cross-bandwidth-dynamics}

\[
\frac{d\gamma_{ij}}{dt} = -\frac{1}{\tau_{\mathrm{dec}}^{(\mathrm{ii})}}(\gamma_{ij} - \gamma_{\mathrm{floor}}) + \xi(t)
\]

with \(\gamma_{\mathrm{floor}}\) the architecturally-supported
steady-state and \(\xi(t)\) stochastic substrate perturbation.

Two regimes:

\begin{itemize}
\tightlist
\item
  \textbf{Sustained-above-threshold:}
  \(\gamma_{\mathrm{floor}} > \Gamma_{\mathrm{min}}\) with margin.
  \(\tau_{(\mathrm{ii})}\) is exponentially long via Kramers-class
  escape from the protected configuration. \emph{SC qubit T₂ regime.}
\item
  \textbf{Near-threshold:}
  \(\gamma_{\mathrm{floor}} \approx \Gamma_{\mathrm{min}}\), by
  architectural design (Josephson junction). Fluctuations frequently
  push below threshold; each crossing produces a global rule
  reconfiguration. \textbf{WKB MQT structure recovered directly from
  DCGT.}
\end{itemize}

\subsubsection{4.3 Condition (iii): Commitment-Injection
Dynamics}\label{condition-iii-commitment-injection-dynamics}

\[
\tau_{(\mathrm{iii})} = \frac{\ln(1/\mathcal{U}_{\mathrm{min}}^{(\mathrm{iii})})}{\Lambda_{\mathrm{env}} + \Lambda_{\mathrm{int}}}.
\]

\(\Lambda_{\mathrm{env}}\) has additive contributions from thermal
radiation, residual-gas collisions, Purcell decay, dielectric noise,
magnetic-flux fluctuations, quasiparticle injection. Bounded below by V1
vacuum kernel residual. \(\Lambda_{\mathrm{int}}\) from intentional +
incidental commitment events.

\subsubsection{4.4 The QC Operating
Window}\label{the-qc-operating-window}

The QC operating window is the substrate time over which all three UR-1
conditions simultaneously hold:

\[
\boxed{\;\tau_{\mathrm{QC}}(\mathcal{S}, \mathcal{E}) = \min\!\left[\tau_{(\mathrm{i})}, \tau_{(\mathrm{ii})}, \tau_{(\mathrm{iii})}\right].\;}
\]

Three structural observations.

\textbf{Different architectures bind on different conditions.} An
architectural class is characterized not by which timescale is longest
but by which is shortest --- i.e., which UR-1 condition is the
architecture's binding constraint.

\textbf{The operating window is a substrate observable.}
\(\tau_{\mathrm{QC}}\) corresponds to what experimentalists call
coherence time, T₁/T₂, fault-tolerance-headroom, or maximum coherent
flight time. Cross-platform measurements are evaluations of
\(\tau_{\mathrm{QC}}\) for different binding conditions.

\textbf{Architectural design is the optimization of \(\min(\cdot)\).}
Improving the longest of three timescales is irrelevant; only
improvement of the binding constraint extends the window. The substrate
provides the structural reason for what platform-development teams
empirically know.

\subsubsection{4.5 Cross-Anchor
Reductions}\label{cross-anchor-reductions}

{\def\LTcaptype{none} % do not increment counter
\begin{longtable}[]{@{}
  >{\raggedright\arraybackslash}p{(\linewidth - 4\tabcolsep) * \real{0.3333}}
  >{\raggedright\arraybackslash}p{(\linewidth - 4\tabcolsep) * \real{0.3333}}
  >{\raggedright\arraybackslash}p{(\linewidth - 4\tabcolsep) * \real{0.3333}}@{}}
\toprule\noalign{}
\begin{minipage}[b]{\linewidth}\raggedright
Phenomenon
\end{minipage} & \begin{minipage}[b]{\linewidth}\raggedright
Binding condition
\end{minipage} & \begin{minipage}[b]{\linewidth}\raggedright
Substrate identification
\end{minipage} \\
\midrule\noalign{}
\endhead
\bottomrule\noalign{}
\endlastfoot
Matter-wave Q-C boundary 140--250 kDa & (i) static failure &
\(\mathcal{M}_{\mathrm{floor}}(M_\star) = \mathcal{M}_{\mathrm{crit}}\) \\
SC qubit T₁ & (iii) & \(\Lambda_{\mathrm{env}}\) Purcell + dielectric +
quasiparticle \\
SC qubit T₂ & (ii) & Kramers escape from protected pathway \\
JJ MQT & (ii) near-threshold & WKB
\(\sim \omega_0\exp[-\alpha\int\sigma\,d\ell]\) \\
Topological-qubit error suppression & (iii) gap-suppressed & Most env
modes don't couple to protected rule \\
Multi-timescale FPM relaxation (Hafezi 2025) & (i) relaxed by redundancy
& Multi-axis pathway expansion \\
Quasiparticle poisoning & multi-condition (i)+(iii) & Single event
raises \(\mathcal{M}\) and \(\Lambda\) \\
\end{longtable}
}

\begin{center}\rule{0.5\linewidth}{0.5pt}\end{center}

\subsection{5. Architectural-Class Taxonomy and
Exhaustiveness}\label{architectural-class-taxonomy-and-exhaustiveness}

Three substrate-allowed protection strategies emerge as the base over
UR-1's three conditions.

\subsubsection{5.1 Class A ---
Engineered-Low-Multiplicity}\label{class-a-engineered-low-multiplicity}

\textbf{Protection target:} condition (i). The architecture commits
structurally to suppressing local multiplicity via lattice symmetry,
engineered ED-bottleneck, deep confinement, or mode engineering.

\textbf{Substrate mechanism.} A bulk superconductor's lattice symmetry
collapses \(\mathcal{M}\) to ≈ 1 (ED-I-01); a Josephson junction's
barrier imposes a sparsity-bottleneck enforcing non-individuation across
the gap (ED-I-23); a deeply trapped ion occupies a low-\(\mathcal{M}\)
region of substrate; a single-photon mode in a linear-optical circuit
has minimum-channel multiplicity (ED-I-13).

\textbf{Suppression of non-target conditions.} (ii) handled by
engineering low-\(\sigma\) corridors between qubit endpoints; (iii) by
isolation engineering (dilution refrigerators, ultra-high vacuum,
low-loss substrates, magnetic shielding).

\textbf{Binding constraint.} Typically (ii) or (iii) --- environmental
coupling channels reach through the engineered isolation. T₁ and T₂ for
SC qubits are the empirical binding.

\textbf{Examples:} Transmon and fluxonium SC qubits; trapped-ion qubits;
gate-model photonic qubits; matter-wave-interferometry molecules.

\subsubsection{5.2 Class B ---
Global-Geometric-Rigidity}\label{class-b-global-geometric-rigidity}

\textbf{Protection target:} condition (ii). The architecture encodes the
participation rule in \emph{global ED-channel geometry} --- topological
invariants that cannot be perturbed by local environmental injection.

\textbf{Substrate mechanism.} Topological phases (Aharonov-Bohm, Berry,
Aharonov-Casher) are global invariants of multiply-connected ED-channels
with nontrivial curvature (ED-I-14). Photonic Chern channels carry
information through edge-states that remain edge-localized under
perturbation by virtue of topological gradient-rigidity (ED-I-12,
ED-I-18 §7). Anyonic platforms exploit ED-I-14's degeneracy-sourced
parameter-space curvature to produce phase shifts depending on holonomy
of parameter loops, not on local rates.

\textbf{Suppression of non-target conditions.} (i) gap-suppressed:
\(\mathcal{M}_{\mathrm{floor}}\) in the protected sector held low by gap
structure. (iii) gap-suppressed: most environmental modes don't couple
to the rule because it's encoded in global geometry;
\(\Lambda_{\mathrm{env}}\) for protected subspace is exponentially
suppressed in \(\Delta_{\mathrm{top}}\).

\textbf{Binding constraint.} Topology-perturbation rate
\(\tau_{\mathrm{gap-stab}}(\mathcal{T})\) --- gap-closing events,
edge-state hybridization with bulk modes, anyonic non-Abelian braiding
errors. Structurally distinct from Class A: limit is set by macroscopic
substrate quality, not local environmental rate.

\textbf{Examples:} Majorana-class topological qubits; photonic Chern
edge-state channels; geometric-phase quantum gates.

\subsubsection{5.3 Class C ---
High-Multiplicity-Redundancy}\label{class-c-high-multiplicity-redundancy}

\textbf{Protection target:} \(\mathcal{U}\) via redundancy of pathways.
Condition (i) is \textbf{relaxed by architectural design} --- high
\(\mathcal{M}\) is permitted because the rule's integrity survives via
parallel-pathway redundancy.

\textbf{Substrate mechanism.} The architecture configures multiple
parallel ED-channels carrying the same participation rule.
Single-pathway failure leaves the rule intact across surviving pathways.
ED-I-18 establishes this for multi-timescale photonic lattices: a
single-ring resonator is a low-\(\mathcal{M}\) object with strict
frequency-phase matching (FPM); a multi-timescale lattice expands
ED-multiplicity along multiple axes and FPM relaxes from a sharp
condition to a region with 100\% device yield, two-octave bandwidth, and
disorder tolerance.

\textbf{Suppression of non-target conditions.} Effective
\(\Gamma_{\mathrm{cross}}\) enhanced by multiple parallel pathways;
effective \(\Lambda\) for protected information distinct from
\(\Lambda\) for any single mode.

\textbf{Binding constraint.} Correlated errors across redundant pathways
exceeding the \emph{correlation budget} \(N_{\mathrm{corr}}\) ---
common-mode noise, shared substrate defects, optical-mode crosstalk.
Saturation is structural: beyond \(N_{\mathrm{corr}}\), redundancy stops
paying.

\textbf{Examples:} Hafezi multi-timescale photonic lattices; bosonic
codes (cat states, GKP); generalized multi-axis architectures.

\subsubsection{5.4 Exhaustiveness
Argument}\label{exhaustiveness-argument}

The three classes exhaust the substrate-level base protection
strategies:

\begin{enumerate}
\def\labelenumi{\arabic{enumi}.}
\tightlist
\item
  UR-1 has three conditions, set by three substrate quantities
  (\(\mathcal{M}, \sigma\)/\(\gamma_{ij}, \Lambda\)).
\item
  Two of these are \emph{fixable structurally} --- \(\mathcal{M}\) via
  local simplification (Class A); \(\sigma\) via global topological
  invariants (Class B). \(\Lambda\) is not directly fixable structurally
  because it is intrinsically a coupling-rate at the boundary between
  protected region and environment; reducing \(\Lambda\) is achieved as
  a \emph{consequence} of the other classes' mechanisms.
\item
  The third strategy is \emph{indirect} --- provide redundancy so
  \(\mathcal{U}\) survives partial failures (Class C).
\item
  No fourth substrate-level base strategy: only three substrate
  quantities, only two structurally fixable, redundancy provides the
  third option.
\end{enumerate}

\subsubsection{5.5 Meta-Architectures: Compositions and
Techniques}\label{meta-architectures-compositions-and-techniques}

Real systems are typically \emph{compositions} of A/B/C, not pure
single-class instances. A fault-tolerant SC quantum computer is Class A
at the physical-qubit level, Class C in the logical encoding (e.g.,
surface code), with active restoring at the logical level
(syndrome-guided correction). Hybrids (SC-photonic interconnects,
topological-SC platforms) are A/B/C subsystems with substrate-geometric
handoff boundaries.

Four techniques considered as Class D candidates and rejected:

\begin{itemize}
\tightlist
\item
  \textbf{Dynamical decoupling.} Reduces effective \(\Lambda\) via
  temporal averaging. \emph{Technique} for Class A.
\item
  \textbf{Reservoir-engineered unresolvedness} (Mirrahimi-Devoret cat
  states). Implements \(A_{S}\) restoring via structured environmental
  coupling. \emph{Technique} for Class A.
\item
  \textbf{Error-correction-as-architecture.} Class C encoding + Class A
  active restoring at logical level. \emph{Meta-architecture composing
  existing classes.}
\item
  \textbf{Hybrid systems.} \emph{Compositions} of A/B/C across
  subsystems.
\end{itemize}

None introduces a fourth substrate-state quantity to fix.

\begin{center}\rule{0.5\linewidth}{0.5pt}\end{center}

\subsection{\texorpdfstring{6. The Multiplicity-Cap Function
\(M\)}{6. The Multiplicity-Cap Function M}}\label{the-multiplicity-cap-function-m}

\subsubsection{6.1 The Unified Function}\label{the-unified-function}

\(M\) is constructed as one substrate object with three
architectural-class projections:

\[
\boxed{\;M(\mathcal{S}, K, \mathcal{E}, \mathcal{O}) = \min\!\left[\tau_{(\mathrm{i})}^{(K, \mathcal{O})}, \tau_{(\mathrm{ii})}^{(K, \mathcal{O})}, \tau_{(\mathrm{iii})}^{(K, \mathcal{O})}\right]\;}
\]

with inputs: - System \(\mathcal{S}\) (protected region's configuration,
endpoints, \(\mathcal{M}_{\mathrm{floor}}\),
\(\gamma_{\mathrm{floor}}\), redundancy count). - Class
\(K \in \{A, B, C\}\). - Environment \(\mathcal{E}\)
(\(\Lambda_{\mathrm{env}}\), perturbation rates, temperature). -
Meta-architecture overlay \(\mathcal{O}\) (dynamical decoupling,
reservoir engineering, error correction, hybrid composition --- any
combination).

Each \(\tau_j^{(K,\mathcal{O})}\) is evaluated using the rate equation
from Section 4 with class-specific protection-mechanism modifiers.

\subsubsection{\texorpdfstring{6.2 Class A Projection
\(M_A\)}{6.2 Class A Projection M\_A}}\label{class-a-projection-m_a}

\[
\tau_{(\mathrm{i})}^{(A)} = \begin{cases}
\dfrac{\mathcal{M}_{\mathrm{crit}} - \mathcal{M}_{\mathrm{floor}}(\mathcal{S})}{\dot{\mathcal{M}}_{\mathrm{eff}}^{(A)}} & \text{if } \mathcal{M}_{\mathrm{floor}}(\mathcal{S}) < \mathcal{M}_{\mathrm{crit}} \\[4pt]
0 & \text{if } \mathcal{M}_{\mathrm{floor}}(\mathcal{S}) \geq \mathcal{M}_{\mathrm{crit}} \quad\text{(static failure)}
\end{cases}
\]

\[
\tau_{(\mathrm{ii})}^{(A)} = \tau_{\mathrm{dec}}^{(\mathrm{ii})}\,\exp\!\left[\frac{(\gamma_{\mathrm{floor}}^{(A)} - \Gamma_{\mathrm{min}})^2}{2\sigma_\xi^2}\right]
\]

For JJ at engineered (ii) boundary:
\(\tau_{(\mathrm{ii})}^{(A,\mathrm{JJ})} \approx \omega_0^{-1}\exp[+\alpha\int_{\mathrm{barrier}}\sigma\,d\ell]\)
--- WKB MQT structure.

\[
\tau_{(\mathrm{iii})}^{(A)} = \frac{\ln(1/\mathcal{U}_{\mathrm{min}})}{\Lambda_{\mathrm{env}}^{(A)}(\mathcal{E}) + \Lambda_{\mathrm{int}}^{(A)}(\mathcal{S})}.
\]

\subsubsection{\texorpdfstring{6.3 Class B Projection
\(M_B\)}{6.3 Class B Projection M\_B}}\label{class-b-projection-m_b}

\[
\tau_{(\mathrm{i})}^{(B)} = \tau_{(\mathrm{i})}^{\mathrm{trivial}}\cdot\exp(\Delta_{\mathrm{top}}/T_{\mathrm{eff}})
\]

\[
\tau_{(\mathrm{ii})}^{(B)} = \tau_{\mathrm{gap-stab}}(\mathcal{T})
\]

\[
\tau_{(\mathrm{iii})}^{(B)} = \frac{\ln(1/\mathcal{U}_{\mathrm{min}})}{\Lambda_{\mathrm{env}}^\mathrm{baseline}\cdot\exp(-\Delta_{\mathrm{top}}/T_{\mathrm{eff}})}
\]

Binding shifted from environmental rate to topology-perturbation rate
\(\tau_{\mathrm{gap-stab}}(\mathcal{T})\), set by macroscopic substrate
quality. For Majorana platforms with \(\Delta_{\mathrm{top}} \sim 0.1\)
meV at \(T_{\mathrm{eff}} \sim 20\) mK: gap factor
\(\sim e^{58} \sim 10^{25}\).

\subsubsection{\texorpdfstring{6.4 Class C Projection
\(M_C\)}{6.4 Class C Projection M\_C}}\label{class-c-projection-m_c}

Three redundancy modifiers apply to the base rates:

\begin{itemize}
\tightlist
\item
  \textbf{Effect 1:}
  \(\mathcal{M}^{\mathrm{eff}}(\mathcal{S}) = \mathcal{M}_{\mathrm{floor}}(\mathcal{S})\cdot g(N)\)
  with \(g(N) \to 0\) as \(N \to \infty\).
\item
  \textbf{Effect 2:}
  \(\gamma_{\mathrm{eff}}^{(C)} = \gamma_{\mathrm{floor}}\cdot h(N)\)
  with \(h(N) > 1\) for \(N > 1\).
\item
  \textbf{Effect 3:}
  \(\Lambda^{\mathrm{eff}}_{(C)}(N) = \Lambda^{\mathrm{single}}\cdot c(N)\)
  with \(c(N) \sim \exp(-N/N_{\mathrm{corr}})\) for
  \(N \ll N_{\mathrm{corr}}\), saturating at
  \(N \gtrsim N_{\mathrm{corr}}\).
\end{itemize}

Form FORCED for all three; specific shapes INHERITED from the
substrate-coupling pattern across \(N\) pathways.

Class C ceiling:
\(\tau_{\mathrm{QC}}^{C,\mathrm{ceil}} = \tau_{\mathrm{QC}}^{\mathrm{single}}\cdot e\)
asymptotically. \textbf{Class C is bounded-redundancy-advantage} --- it
cannot exceed single-pathway ceilings without bound; its role is
extending the useful operating window before saturation and mitigating
disorder/fabrication variations.

\subsubsection{6.5 Meta-Architectural
Overlays}\label{meta-architectural-overlays}

Meta-architectures appear as boundary-condition modifiers, not new
projections:

\begin{itemize}
\tightlist
\item
  \textbf{Dynamical decoupling.} Replaces
  \(\Lambda_{\mathrm{env}}^{(K)}\) with temporally-averaged
  \(\Lambda_{\mathrm{env}}^{\mathrm{DD}}(\omega_{\mathrm{decoup}})\).
\item
  \textbf{Reservoir engineering.} Increases \(A_{S}\) in the
  multiplicity dynamics.
\item
  \textbf{Error correction.} Recursive overlay:
  \(\tau_{\mathrm{QC}}^{\mathrm{logical}}\) derived from
  \(\tau_{\mathrm{QC}}^{\mathrm{phys}}\) via code distance + threshold.
\item
  \textbf{Hybrid composition.} \(\min\) over subsystems' projections +
  handoff losses.
\end{itemize}

\begin{center}\rule{0.5\linewidth}{0.5pt}\end{center}

\subsection{7. Cross-Platform Unification and Predictive
Content}\label{cross-platform-unification-and-predictive-content}

\subsubsection{7.1 The Matter-Wave / Qubit-System
Identity}\label{the-matter-wave-qubit-system-identity}

The matter-wave Q-C boundary at 140--250 kDa and qubit-system
multiplicity walls are \textbf{the same boundary projected onto two
platforms}. Both arise from
\(\mathcal{M}_{\mathrm{floor}}^{(A)}(\mathcal{S}) = \mathcal{M}_{\mathrm{crit}}\)
--- the static-failure branch of the Class A \(\tau_{(\mathrm{i})}\)
projection.

{\def\LTcaptype{none} % do not increment counter
\begin{longtable}[]{@{}
  >{\raggedright\arraybackslash}p{(\linewidth - 6\tabcolsep) * \real{0.2500}}
  >{\raggedright\arraybackslash}p{(\linewidth - 6\tabcolsep) * \real{0.2500}}
  >{\raggedright\arraybackslash}p{(\linewidth - 6\tabcolsep) * \real{0.2500}}
  >{\raggedright\arraybackslash}p{(\linewidth - 6\tabcolsep) * \real{0.2500}}@{}}
\toprule\noalign{}
\begin{minipage}[b]{\linewidth}\raggedright
Platform
\end{minipage} & \begin{minipage}[b]{\linewidth}\raggedright
\(\mathcal{S}\)
\end{minipage} & \begin{minipage}[b]{\linewidth}\raggedright
\(f_{\mathrm{sys}}^{(A)}(\mathcal{S})\) scaling parameter
\end{minipage} & \begin{minipage}[b]{\linewidth}\raggedright
Q-C boundary
\end{minipage} \\
\midrule\noalign{}
\endhead
\bottomrule\noalign{}
\endlastfoot
Matter-wave & molecule of mass \(M_{\mathrm{mol}}\) & molecular
internal-DOF activation \(f_{\mathrm{int}}(M_{\mathrm{mol}})\) &
140--250 kDa (empirical) \\
SC qubit array & \(N\)-qubit register & cross-qubit pathway count
\(f_{\mathrm{xy}}(N_{\mathrm{qubits}}, \mathrm{architecture})\) &
\(N_\star\) qubits (INHERITED) \\
Trapped-ion array & \(N\)-ion crystal & inter-ion + motional-mode
coupling & analogous \(N_\star\) \\
Photonic gate-model & photonic-mode register & mode-coupling pathway
count & analogous \(N_\star\) \\
\end{longtable}
}

\textbf{Sharp cross-platform claim:} the matter-wave boundary value
(140--250 kDa), calibrated against \(\mathcal{M}_{\mathrm{crit}}\),
fixes the substrate-determined cap on Class A platforms across all
instantiations. Sharper measurement of the matter-wave boundary directly
constrains the system-multiplicity wall location for Class A qubit
platforms via the substrate-level identity. Qubit-platform \(N_\star\)
values are not independent free parameters --- they are determined (up
to architecture-specific calibration) by the matter-wave anchor.

This is the kind of cross-platform consistency that distinguishes
substrate-level frameworks from phenomenological assemblies.

\subsubsection{7.2 Class-Specific Scaling Laws and
Ceilings}\label{class-specific-scaling-laws-and-ceilings}

\textbf{Class A.} System-multiplicity wall: as
\(f_{\mathrm{sys}}^{(A)}(\mathcal{S}) \to \mathcal{M}_{\mathrm{crit}}\),
\(\tau_{(\mathrm{i})}^{(A)} \to 0\) regardless of environmental
engineering. Improving isolation, dynamical decoupling, or any
environmental-side technique cannot push past this wall. \emph{Form
FORCED; wall location INHERITED.}

\textbf{Class B.} Exponential coherence enhancement by
\(\exp(\Delta_{\mathrm{top}}/T_{\mathrm{eff}})\) over Class A; binding
shifted to topology-perturbation rate. Engineering effort that improves
\(\Delta_{\mathrm{top}}\) or topology stability produces exponential
improvements; effort that improves isolation in the Class A sense
produces diminishing returns. \emph{Form FORCED; max gap and stability
INHERITED.}

\textbf{Class C.} Correlation-budget plateau: exponential improvement
for \(N \ll N_{\mathrm{corr}}\), saturation at
\(\sim e\cdot\tau^{\mathrm{single}}\) for
\(N \gtrsim N_{\mathrm{corr}}\). Bounded-redundancy-advantage.
\emph{Form FORCED; \(N_{\mathrm{corr}}\) INHERITED.}

\subsubsection{7.3 Cross-Class Architectural
Transitions}\label{cross-class-architectural-transitions}

\textbf{Class A → Class C is mandatory at the wall.} When
\(f_{\mathrm{sys}}^{(A)}(\mathcal{S}) \to \mathcal{M}_{\mathrm{crit}}\),
the only path past the wall is to \emph{accept} high physical-level
\(\mathcal{M}\) and use redundancy to maintain logical-level
\(\mathcal{U}\). Sharp prediction: a SC-only architecture without
logical encoding will plateau in effective
\(\tau_{\mathrm{QC}}^{\mathrm{logical}}\) before reaching the
engineering frontier of physical-qubit count, even with continued T₁/T₂
improvements.

\textbf{Class B overtakes Class A at
\(\Delta_{\mathrm{top}}/T_{\mathrm{eff}} \gtrsim 10\)--20.} For
\(T_{\mathrm{eff}} \sim 20\) mK, gaps \(\gtrsim 0.5\)--1 meV with stable
topology on millisecond timescales would push Class B past current Class
A coherence times by orders of magnitude.

\textbf{Class C saturation at \(N \sim N_{\mathrm{corr}}\).}
Hafezi-class platforms scaling beyond \textasciitilde3--5 timescale axes
will show diminishing returns. Bosonic codes scaling beyond a
code-distance threshold will exhibit saturation due to correlated
photon-loss.

\subsubsection{7.4 Falsifiable Predictions and Near-Term Experimental
Tests}\label{falsifiable-predictions-and-near-term-experimental-tests}

\textbf{Test 1 --- Coherence-scaling experiments (Class A wall).}
Measure \(\tau_{\mathrm{QC}}^{\mathrm{logical}}\) in scaled SC qubit
systems vs \(N_{\mathrm{qubits}}\) at fixed code distance, with
continued T₁/T₂ engineering. Predicted plateau as
\(f_{\mathrm{xy}}(N_{\mathrm{qubits}})\) approaches
\(\mathcal{M}_{\mathrm{crit}}\). Continued unbounded improvement past
the projected wall would refute.

\textbf{Test 2 --- Multiplicity-growth probes.} Mass-matched molecular
isomers with different DOF structure for matter waves; cross-qubit
pathway tomography for qubit systems. Tests whether contrast/coherence
scales with substrate-level \(\mathcal{M}^{\mathrm{system}}\), not
directly with mass/qubit-count.

\textbf{Test 3 --- Topology-perturbation thresholds (Class B).} Measure
\(\tau_{\mathrm{QC}}\) in topological qubits as a function of
\(\Delta_{\mathrm{top}}\) and temperature. Predicted exponential
dependence on \(\Delta_{\mathrm{top}}/T_{\mathrm{eff}}\). Polynomial
dependence would refute.

\textbf{Test 4 --- Redundancy-saturation signatures (Class C).} Extend
Hafezi-class to higher axis counts; scale bosonic-code distance.
Predicted plateau emerges. Continued monotonic improvement past
projected saturation would refute.

\subsubsection{7.5 Decisive Falsifier of the
Framework}\label{decisive-falsifier-of-the-framework}

The arc's central falsifier: identification of a platform whose
coherence behavior is consistent with no class assignment in \{A, B,
C\}, cannot be explained as meta-architectural composition, and fails to
match \(M\) under any substrate-allowed parameter calibration. Such a
platform forces introduction of a Class D (refuting Section 5's
exhaustiveness) or replacement of the framework. Conversely, continued
cross-anchor consistency over the next decade of platform development is
significant cross-platform evidence for the substrate framework.

\begin{center}\rule{0.5\linewidth}{0.5pt}\end{center}

\subsection{8. FORCED / INHERITED / OPEN
Classification}\label{forced-inherited-open-classification}

\subsubsection{FORCED (substrate-level
form)}\label{forced-substrate-level-form}

\begin{itemize}
\tightlist
\item
  UR-1 as necessary-and-sufficient characterization of the unresolved
  regime.
\item
  Three-factor product structure of \(\mathcal{U}\).
\item
  \(\min(\cdot)\) structure of \(\tau_{\mathrm{QC}}\).
\item
  Substrate-level dynamics of multiplicity, cross-bandwidth, and
  commitment-rate.
\item
  Static-failure branch of \(\tau_{(\mathrm{i})}\) at
  \(\mathcal{M}_{\mathrm{floor}} = \mathcal{M}_{\mathrm{crit}}\).
\item
  WKB-form exponential structure of MQT recovered from DCGT.
\item
  Three-class architectural exhaustiveness.
\item
  Meta-architectural overlay structure as modifiers, not new classes.
\item
  Multiplicity-cap function \(M\) as one substrate object with three
  projections.
\item
  Exponential gap-suppression in Class B by
  \(\exp(\Delta_{\mathrm{top}}/T_{\mathrm{eff}})\).
\item
  Correlation-budget plateau in Class C; bounded redundancy advantage.
\item
  Cross-platform unification: matter-wave Q-C boundary and qubit-system
  walls are one phenomenon.
\item
  Three cross-class architectural transitions.
\item
  Cross-domain echo with BH-2 (\(\Gamma_{\mathrm{cross}}\) collapse
  mechanism).
\end{itemize}

\subsubsection{INHERITED (value-level)}\label{inherited-value-level}

\begin{itemize}
\tightlist
\item
  \(\mathcal{M}_{\mathrm{crit}}\), \(\Gamma_{\mathrm{min}}\),
  \(\Lambda_{\mathrm{V1}}\) --- substrate constants.
\item
  Specific functional shapes of \(\mu, \kappa, g, h, c\).
\item
  Matter-wave Q-C boundary location (140--250 kDa).
\item
  Architecture-specific \(f_{\mathrm{sys}}^{(A)}\) scaling functions.
\item
  Specific values of \(N_\star\) for any qubit platform.
\item
  Maximum stable topological gap \(\Delta_{\mathrm{top}}^{\max}\),
  minimum substrate-equivalent temperature \(T_{\mathrm{eff,\min}}\).
\item
  Specific \(N_{\mathrm{corr}}\) for any Class C platform.
\item
  Specific cross-class transition thresholds (numerical).
\end{itemize}

\subsubsection{OPEN (future arcs)}\label{open-future-arcs}

\begin{itemize}
\tightlist
\item
  \textbf{O-QC-1:} Closed-form \(\mathcal{M}_{\mathrm{crit}}\) from V1
  kernel + ED-I-01 substrate constants. Structurally similar in
  difficulty to closed-form \(\log g\) (BH-5).
\item
  \textbf{O-QC-2:} Architecture-to-platform calibration program --- map
  \(\mathcal{M}_{\mathrm{floor}}^{(A)}(\mathcal{S})\) for canonical
  platforms; predict \(N_\star\) values.
\item
  \textbf{O-QC-3:} Closed-form \(g, h, c(N)\) for Class C from extended
  DCGT.
\item
  \textbf{O-QC-4:} Topology-stability theorem for Class B;
  substrate-level account of \(\tau_{\mathrm{gap-stab}}(\mathcal{T})\).
\item
  \textbf{O-QC-5:} Surface-code logical-qubit recursive-overlay
  derivation.
\item
  \textbf{O-QC-6:} Hybrid-architecture handoff dynamics.
\end{itemize}

\begin{center}\rule{0.5\linewidth}{0.5pt}\end{center}

\subsection{9. Discussion and
Implications}\label{discussion-and-implications}

\subsubsection{9.1 Relation to Standard QC
Theory}\label{relation-to-standard-qc-theory}

ED's substrate-level account does not predict that current
quantum-computing experiments will give different numerical answers than
standard QC theory predicts. SC qubit T₁/T₂ measurements will continue
to look as they do; matter-wave interferometry will continue to lose
coherence at the molecular masses it does; topological-qubit progress
will continue to be paced by topology engineering. What changes is the
\emph{story underneath}: the question of why a particular platform hits
a particular wall has a substrate-level structural answer, and the
cross-platform identity of those walls is now visible.

\subsubsection{9.2 Why ED Doesn't Solve QC Engineering --- and That Is
the
Point}\label{why-ed-doesnt-solve-qc-engineering-and-that-is-the-point}

This paper does not deliver a constructive recipe for building a
fault-tolerant quantum computer. The arc is structural, not engineering.
What it does deliver is a framework in which engineering questions
become substrate-resolvable:

\begin{itemize}
\tightlist
\item
  \emph{Is there a hard wall on SC qubit scaling?} Yes, at the
  system-multiplicity boundary; same boundary as the matter-wave Q-C
  boundary.
\item
  \emph{Will adding more redundancy axes keep paying?} No, past the
  correlation budget; bounded-redundancy-advantage is structural.
\item
  \emph{Can topological qubits exceed SC coherence by a lot?} Yes in
  principle, by exponential factors; bound is topology engineering, not
  isolation.
\item
  \emph{Is error correction a different kind of solution from
  physical-qubit engineering?} No, it is a meta-architectural
  composition of the existing classes at logical level.
\end{itemize}

These answers are substrate-grounded, not phenomenological inferences
from existing data. They make sharp, testable predictions about which
engineering directions will continue to pay and which will not.

\subsubsection{9.3 The Cross-Domain Echo with Black-Hole
Architecture}\label{the-cross-domain-echo-with-black-hole-architecture}

The same substrate mechanism --- cross-bandwidth
\(\Gamma_{\mathrm{cross}}\) collapse via DCGT-derived
\(\exp[-\alpha\sigma]\) structure --- underlies black-hole horizon
formation (Arc BH, BH-2) and condition (ii) failure in QC. Two faces of
one substrate identity, applied at scales separated by \textasciitilde50
orders of magnitude in physical length. This is the cross-domain
consistency a substrate-level framework should produce; phenomenological
frameworks cannot deliver it.

The structural lesson: the substrate doesn't care whether the
gradient-collapse is happening at a black-hole event horizon or at a
Josephson-junction barrier. The same DCGT machinery describes both.
Different empirical phenomena, different physical scales, one substrate
mechanism.

\subsubsection{9.4 The QC Scaling Debates,
Reframed}\label{the-qc-scaling-debates-reframed}

A recurring tension in quantum-computing research has been the debate
between optimists (continued scaling will deliver fault tolerance
through engineering) and skeptics (a fundamental limit prevents
practical-scale QC). The ED account gives a structural framework that
resolves the debate's framing without taking sides on whether
practical-scale QC is achievable.

There \emph{is} a fundamental limit --- the substrate-level multiplicity
wall --- but it is not where skeptics typically place it. The wall is
platform-architecture-dependent. Class A platforms hit a wall in
physical-qubit count; Class B platforms can in principle exceed it by
exponential factors at the cost of much harder engineering; Class C
platforms hit a different wall (correlation-budget saturation) at a much
lower scale. The optimists' continued-scaling expectation is correct
\emph{for Class A platforms below the wall}, incorrect \emph{for Class A
above the wall without composition}, and irrelevant \emph{for Class B/C}
which face structurally different binding constraints.

The framework dissolves the binary debate by structurally locating the
limits --- different limits, on different timescales, for different
architectures --- rather than resolving it within a single-platform
engineering frame.

\subsubsection{9.5 What This Changes for the Engineering
Frontier}\label{what-this-changes-for-the-engineering-frontier}

Three near-term implications for platform development:

\begin{itemize}
\tightlist
\item
  \textbf{For SC programs:} the system-multiplicity wall is a real
  structural feature; engineering effort beyond the wall must include
  composition with Class C (logical encoding) --- pure-Class-A scaling
  will plateau.
\item
  \textbf{For topological qubit programs:} the binding is \emph{topology
  engineering}, not isolation engineering; effort to improve
  \(\Delta_{\mathrm{top}}\) and gap stability is the lever, not effort
  to improve coupling-suppression.
\item
  \textbf{For multi-axis / bosonic-code programs:} the
  bounded-redundancy-advantage is structural; sub-linear scaling past
  saturation is predicted, and this should be measured directly to
  validate or falsify the framework.
\end{itemize}

\subsubsection{9.6 Implications for the Quantum-Computing
Roadmap}\label{implications-for-the-quantum-computing-roadmap}

The multiplicity-cap function \(M\) implies that the dominant scaling
limits of current quantum-computing platforms are \textbf{structural
rather than technological}. The walls that platforms encounter at scale
are not engineering accidents to be removed by improved isolation,
fabrication, or error correction within their current architectural
class --- they are substrate-determined ceilings whose locations are set
by substrate constants (\(\mathcal{M}_{\mathrm{crit}}\),
\(\Gamma_{\mathrm{min}}\), \(N_{\mathrm{corr}}\),
\(\Delta_{\mathrm{top}}^{\max}\)) and whose existence is FORCED by the
substrate machinery.

Five claims follow directly from \(M\) and have direct roadmap
consequences:

\begin{enumerate}
\def\labelenumi{\arabic{enumi}.}
\item
  \textbf{The matter-wave Q-C boundary and qubit-platform multiplicity
  walls are the same substrate phenomenon.} Sharper measurement of the
  matter-wave boundary at 140--250 kDa directly constrains the
  system-multiplicity wall location for Class A qubit platforms.
  Progress in matter-wave interferometry and progress in
  superconducting-qubit scaling are not independent research programs;
  they probe the same substrate quantity from two platforms.
\item
  \textbf{Only Class B can outrun the Class A multiplicity wall by
  structural margin.} Class A platforms (superconducting qubits, trapped
  ions, photonic gate-model) face an exponential collapse of
  \(\tau_{(\mathrm{i})}\) as their effective system-multiplicity
  approaches \(\mathcal{M}_{\mathrm{crit}}\). No amount of isolation
  engineering, dynamical decoupling, or reservoir engineering pushes
  past this wall --- these are techniques that extend Class A's
  environmental-binding timescales, not its multiplicity-binding
  timescale. Pushing past the wall requires either Class B's topological
  gradient-rigidity (binding shifted to topology engineering) or
  composition with Class C (logical encoding via redundancy, which is
  what error correction provides).
\item
  \textbf{Redundancy saturates.} Class C platforms exhibit a
  correlation-budget plateau at \(N \sim N_{\mathrm{corr}}\). Beyond
  this, additional redundancy stops paying --- multi-axis photonic
  platforms above \textasciitilde3--5 axes and bosonic codes above a
  code-distance threshold will exhibit diminishing returns due to
  correlated noise channels. The redundancy advantage is bounded; Class
  C is not an unlimited-scaling strategy.
\item
  \textbf{The fault-tolerant-QC roadmap has a structural shape.}
  Surface-code-based fault-tolerant quantum computers are
  meta-architectural compositions of Class A (physical-qubit level) with
  Class C (logical encoding). Their scaling is governed by the
  substrate-level identity of their constituent layers: physical-layer
  Class A walls bound the substrate of the composition, and the
  logical-layer Class C correlation-budget bounds the gain from
  encoding. Strategies that aim to push the composition past \emph{both}
  substrate-level bounds simultaneously must include a Class B
  subsystem, since Class A and Class C alone cannot.
\item
  \textbf{Long-term scaling strategy centers on topology-stabilized
  architectures.} Class B is the only architectural class whose ceiling
  is set by \emph{macroscopic substrate quality} (gap engineering,
  sample homogeneity) rather than by environmental coupling rates. Its
  exponential scaling advantage over Class A --- by
  \(\exp(\Delta_{\mathrm{top}}/T_{\mathrm{eff}})\) --- makes it the
  structurally distinct path past the matter-wave-class wall. The cost
  is much harder substrate engineering; the structural reward is a
  coherence ceiling not bounded by isolation-class techniques.
\end{enumerate}

This unification suggests a re-evaluation of long-term scaling
strategies in the quantum-computing field. The emphasis on isolation
engineering and error-correcting overhead has historically been
platform-specific and engineering-driven; the substrate-level account
argues that this emphasis is correct \emph{within Class A} but bounded
by a structural wall whose location is empirically constrained by the
matter-wave anchor. The path past the wall is either
composition-with-Class-C (already the standard fault-tolerant approach,
but bounded by correlation-budget saturation) or Class B
(topology-stabilized architectures, with their structurally distinct
binding mode).

The framework neither predicts that practical-scale quantum computing is
achievable nor that it is impossible. What it does predict is that
\emph{which substrate-level wall is binding at what scale} determines
what engineering effort will pay and what will not --- and the
cross-platform identity of these walls means that empirical progress in
any one domain (matter-wave interferometry, SC scaling,
topological-qubit gap stability, multi-axis-redundancy saturation)
directly informs the substrate-level constraints on every other domain.

\subsubsection{9.7 Roadmap for Numerical
Calibration}\label{roadmap-for-numerical-calibration}

The structural results presented here --- UR-1, the three-class
taxonomy, the multiplicity-cap function \(M\), and the cross-platform
unification --- are complete in form. The numerical values that fix the
locations of the substrate-level ceilings, however, are INHERITED rather
than closed-form derived in this paper, and they require empirical
calibration before quantitative platform-specific predictions become
sharp. These calibrations are engineering-accessible measurements, not
theoretical unknowns; the paper's role is to provide the structural form
into which the calibrated values are inserted.

Five quantities, five calibration paths:

\begin{itemize}
\tightlist
\item
  \textbf{\(\mathcal{M}_{\mathrm{crit}}\)} --- calibrated by matter-wave
  mass-cap refinement (sharper measurement of the 140--250 kDa Q-C
  boundary using mass-matched isomers with controlled internal-DOF
  structure) combined with qubit-platform multiplicity tomography
  (cross-talk and leakage characterization that maps physical
  \(\mathcal{M}_{\mathrm{floor}}^{(A)}\)). The two anchors converge on
  the same substrate constant via the cross-platform identity
  established in §7.
\item
  \textbf{\(N_{\mathrm{corr}}\)} --- calibrated by redundancy-saturation
  experiments in Class C platforms: extending Hafezi-class
  multi-timescale photonic lattices to higher axis counts and measuring
  where FPM-relaxation, harmonic-bandwidth, and yield improvements
  plateau; analogously, scaling bosonic-code distance and measuring
  logical-qubit lifetime saturation.
\item
  \textbf{\(\Delta_{\mathrm{top}}^{\max}\) and
  \(T_{\mathrm{eff,\min}}\)} --- calibrated by gap-stability
  measurements in Class B platforms: gap spectroscopy in
  topological-qubit candidate materials combined with
  quasiparticle-poisoning and edge-state-hybridization rate measurements
  at varying temperature, mapping the achievable region of the
  \((\Delta_{\mathrm{top}}, T_{\mathrm{eff}})\) plane.
\item
  \textbf{Architecture-specific \(f_{\mathrm{sys}}^{(A)}(\mathcal{S})\)}
  --- calibrated by cross-platform comparison of how effective
  system-multiplicity scales with system size for canonical Class A
  platforms (transmon, fluxonium, trapped ion, photonic gate-model).
  Each architecture has its own scaling, and the cross-architecture
  comparison is itself informative for the substrate-coupling pattern.
\item
  \textbf{Platform-specific \(N_\star\)} --- calibrated by
  system-multiplicity scaling experiments in Class A qubit platforms:
  measuring effective \(\tau_{\mathrm{QC}}^{\mathrm{logical}}\) as
  physical-qubit count rises at fixed code distance and locating the
  plateau predicted in §7.4.
\end{itemize}

These five calibration paths are mutually consistent with the
cross-platform identity: the matter-wave anchor and the qubit-system
anchor must yield consistent values of \(\mathcal{M}_{\mathrm{crit}}\),
and the four primary empirical anchors must lie at substrate-determined
locations on \(M\) across the calibrated parameter space. Calibration is
downstream work; the structural framework is delivered.

\subsubsection{9.8 Topology-Stability as a Downstream
Theorem}\label{topology-stability-as-a-downstream-theorem}

Class B's ceiling is set by the topology-perturbation rate
\(\tau_{\mathrm{gap-stab}}(\mathcal{T})\), which the present arc names
as a structural quantity but does not derive at platform-level
specificity. A full \textbf{topology-stability theorem} would formalize
how global ED-channel geometry resists local substrate perturbation
across canonical topological structures (Majorana, Fibonacci anyons,
Chern bands), expressing \(\tau_{\mathrm{gap-stab}}(\mathcal{T})\) as a
substrate-level functional of the topological invariant \(\mathcal{T}\)
and the substrate's local fluctuation spectrum.

Such a theorem is a natural follow-on arc (priority item O-QC-4) and is
explicitly not solved in this paper. The architectural role of Class B
is complete: its protection mechanism is identified (rigidity of
ED-channel geometry under topological invariants, ED-I-14 + ED-I-18 §7),
its binding constraint is named (topology-perturbation rather than
environmental rate), and its substrate-level scaling is established
(\(\exp(\Delta_{\mathrm{top}}/T_{\mathrm{eff}})\) enhancement over Class
A). The topology-stability theorem provides the platform-level bound on
\(\tau_{\mathrm{gap-stab}}(\mathcal{T})\) that converts the structural
Class B ceiling into a closed-form numerical prediction for any specific
topological architecture. The structural framework is in hand; the
theorem that calibrates it remains downstream work, requiring the
platform-specific geometric analysis that ED-I-14's interpretive content
sketches but does not formalize.

\begin{center}\rule{0.5\linewidth}{0.5pt}\end{center}

\subsection{10. Appendices}\label{appendices}

\subsubsection{Appendix A --- UR-1 Proof
Sketch}\label{appendix-a-ur-1-proof-sketch}

\textbf{Sufficient direction.} Assume UR-1's three conditions (i), (ii),
(iii) hold. Then:

\begin{itemize}
\tightlist
\item
  By (i),
  \(\mu(\mathcal{M}_i/\mathcal{M}_{\mathrm{crit}}) \geq \mu(\mu^{-1}(\mathcal{U}_{\mathrm{min}}^{1/n})) = \mathcal{U}_{\mathrm{min}}^{1/n}\)
  for each endpoint, so the multiplicity-headroom product satisfies
  \(\prod_i \mu \geq \mathcal{U}_{\mathrm{min}}\).
\item
  By (ii), an analogous bound on the connectivity product gives
  \(\prod_{(i,j) \in \mathcal{R}} \kappa \geq \mathcal{U}_{\mathrm{min}}\).
\item
  By (iii),
  \(\exp[-\int_0^t \Lambda dt'] \geq \mathcal{U}_{\mathrm{min}} \cdot e^{C_{(\mathrm{i})} + C_{(\mathrm{ii})}}\).
\end{itemize}

Multiplying the three factors and accounting for cross-condition
deficits:
\(\mathcal{U}(\mathcal{S},t) \geq \mathcal{U}_{\mathrm{min}}\). ∎

\textbf{Necessary direction.} Suppose
\(\mathcal{U}(\mathcal{S},t) \geq \mathcal{U}_{\mathrm{min}}\) but
condition (i) fails at some endpoint \(e_k\):
\(\mathcal{M}_k > \mathcal{M}_{\mathrm{crit}}^{(\mu^{-1}(\mathcal{U}_{\mathrm{min}}^{1/n}))}\).
Then
\(\mu(\mathcal{M}_k/\mathcal{M}_{\mathrm{crit}}) < \mathcal{U}_{\mathrm{min}}^{1/n}\),
so the multiplicity-headroom product is below
\(\mathcal{U}_{\mathrm{min}}\). Even with the other two factors at their
maximum (1 each), \(\mathcal{U} < \mathcal{U}_{\mathrm{min}}\),
contradiction. The same argument applies for (ii) and (iii). ∎

The full closed-form proof requires the specific functional shapes of
\(\mu\) and \(\kappa\), which are INHERITED. The structural argument is
form-FORCED.

\subsubsection{Appendix B --- DCGT-WKB Recovery of
MQT}\label{appendix-b-dcgt-wkb-recovery-of-mqt}

The Josephson-junction MQT rate inherits its WKB exponential structure
directly from DCGT's cross-bandwidth formula. For a JJ engineered at the
(ii) boundary with
\(\gamma_{\mathrm{floor}} \to \Gamma_{\mathrm{min}}\), the dominant
contribution to \(\tau_{(\mathrm{ii})}^{(A,\mathrm{JJ})}\) comes from
below-threshold excursions of \(\gamma_{LR}\). Each excursion produces a
global rule reconfiguration (per ED-I-29). The rate of such excursions
is set by the DCGT \(\kappa\)-function evaluated at the engineered
\(\sigma\)-profile of the barrier:

\[
\tau_{\mathrm{MQT}}^{-1} = \omega_0\,\exp\!\left[-\alpha\!\int_{\mathrm{barrier}}\sigma(\mathbf{x})\,d\ell\right]
\]

with \(\omega_0\) a substrate-attempt frequency (INHERITED) and
\(\alpha\) the DCGT prefactor. The integrand
\(\sigma = |\nabla\rho|\ell_P^2/\rho_{\mathrm{local}}\) is evaluated
across the barrier; in standard JJ language, this maps to barrier height
and width, recovering the standard WKB form.

\subsubsection{Appendix C --- Cross-Anchor Consistency
Table}\label{appendix-c-cross-anchor-consistency-table}

{\def\LTcaptype{none} % do not increment counter
\begin{longtable}[]{@{}
  >{\raggedright\arraybackslash}p{(\linewidth - 10\tabcolsep) * \real{0.1667}}
  >{\raggedright\arraybackslash}p{(\linewidth - 10\tabcolsep) * \real{0.1667}}
  >{\raggedright\arraybackslash}p{(\linewidth - 10\tabcolsep) * \real{0.1667}}
  >{\raggedright\arraybackslash}p{(\linewidth - 10\tabcolsep) * \real{0.1667}}
  >{\raggedright\arraybackslash}p{(\linewidth - 10\tabcolsep) * \real{0.1667}}
  >{\raggedright\arraybackslash}p{(\linewidth - 10\tabcolsep) * \real{0.1667}}@{}}
\toprule\noalign{}
\begin{minipage}[b]{\linewidth}\raggedright
Anchor
\end{minipage} & \begin{minipage}[b]{\linewidth}\raggedright
Class
\end{minipage} & \begin{minipage}[b]{\linewidth}\raggedright
Binding condition
\end{minipage} & \begin{minipage}[b]{\linewidth}\raggedright
Substrate identification
\end{minipage} & \begin{minipage}[b]{\linewidth}\raggedright
Form FORCED
\end{minipage} & \begin{minipage}[b]{\linewidth}\raggedright
Value INHERITED
\end{minipage} \\
\midrule\noalign{}
\endhead
\bottomrule\noalign{}
\endlastfoot
Bulk SC \(T_c\) (ED-I-01) & A & (iii) & thermal
\(\Lambda_{\mathrm{env}}\) crosses threshold & yes & \(T_c\) value \\
JJ MQT (Devoret-Martinis-Clarke) & A & (ii) near-threshold & WKB from
DCGT & yes & prefactor, barrier \(\sigma\) \\
SC qubit T₁ & A & (iii) & Purcell + dielectric + QP & yes & rates \\
SC qubit T₂ & A & (ii) & Kramers escape & yes &
\(\tau_{\mathrm{dec}}^{(\mathrm{ii})}\) \\
Matter-wave Q-C 140--250 kDa & A & (i) static &
\(\mathcal{M}_{\mathrm{floor}} = \mathcal{M}_{\mathrm{crit}}\) & yes &
\(M_\star\) value \\
Topological-gap suppression & B & (iii)/gap-stab &
\(\Lambda \cdot \exp(-\Delta/T_{\mathrm{eff}})\) & yes & \(\Delta\),
\(\tau_{\mathrm{gap-stab}}\) \\
Photonic Chern channels & B & (ii) edge-state rigidity & global
ED-channel invariant & yes & sample quality \\
Hafezi multi-timescale & C & (i) relaxed & \(g(N), h(N), c(N)\) & yes &
\(N_{\mathrm{corr}}\) \\
Bosonic codes (cat, GKP) & C & redundancy budget & photon-number
redundancy & yes & code parameters \\
Surface-code logical & Meta: A+C & logical binding & recursive overlay &
yes & code distance \\
Reservoir-engineered cat & A (overlay) & (i) with active restoring &
\(A_{S}\) via env coupling & yes & engineering details \\
Dynamically-decoupled SC & A (overlay) & (iii) with reduced
\(\Lambda^{\mathrm{eff}}\) & temporal averaging & yes & pulse
parameters \\
\end{longtable}
}

\subsubsection{Appendix D --- Comparison with Standard QC
Frameworks}\label{appendix-d-comparison-with-standard-qc-frameworks}

{\def\LTcaptype{none} % do not increment counter
\begin{longtable}[]{@{}
  >{\raggedright\arraybackslash}p{(\linewidth - 4\tabcolsep) * \real{0.3333}}
  >{\raggedright\arraybackslash}p{(\linewidth - 4\tabcolsep) * \real{0.3333}}
  >{\raggedright\arraybackslash}p{(\linewidth - 4\tabcolsep) * \real{0.3333}}@{}}
\toprule\noalign{}
\begin{minipage}[b]{\linewidth}\raggedright
Feature
\end{minipage} & \begin{minipage}[b]{\linewidth}\raggedright
Standard QC framing
\end{minipage} & \begin{minipage}[b]{\linewidth}\raggedright
ED substrate framing
\end{minipage} \\
\midrule\noalign{}
\endhead
\bottomrule\noalign{}
\endlastfoot
Coherence & State property of wavefunction & Integrity of unresolved
participation rule across endpoints \\
Decoherence & Coupling of state to environment & P11 commitment events
firing at endpoints \\
Engineering target & Reduce coupling-rate to environment & Hold
substrate region in low-\(\mathcal{M}\) + high-\(\gamma_{ij}\) +
low-\(\Lambda\) \\
Coherence-time & T₁ / T₂ as state-decay rates &
\(\tau_{\mathrm{QC}} = \min(\tau_{(\mathrm{i})}, \tau_{(\mathrm{ii})}, \tau_{(\mathrm{iii})})\)
over substrate-condition crossings \\
Architectural classes & platform-specific (SC, ion, photonic, topo) &
substrate-level (A: low-\(\mathcal{M}\); B: rigid-\(\sigma\); C:
redundancy) \\
Cross-platform unification & not addressed & matter-wave Q-C boundary ↔
qubit-system \(N_\star\) on one \(M\) function \\
Hard limit on QC & open question / debated & substrate-determined wall
in Class A; gap-stability in Class B; correlation budget in Class C \\
Error correction & logical encoding for fault tolerance & Class A +
Class C meta-architectural composition \\
MQT rate & WKB postulated from QM & WKB derived from DCGT
cross-bandwidth structure \\
\end{longtable}
}

\begin{center}\rule{0.5\linewidth}{0.5pt}\end{center}

\subsection{11. Metadata}\label{metadata}

\textbf{Keywords.} Event Density; Quantum Computation; Substrate
Physics; Unresolved Regime; UR-1; Multiplicity-Cap Function;
Architectural Class Taxonomy; Matter-Wave Q-C Boundary; Josephson
Junction MQT; Topological Qubits; Multi-Timescale Photonics; Form-FORCED
/ Value-INHERITED Methodology.

\textbf{Version.} 1.0 --- May 2026.

\textbf{Arc Q-COMPUTE references.} Seven memos in
\texttt{theory/Quantum\_Computing/}: Arc\_QC\_1\_Opening;
Arc\_QC\_2\_UR1\_Unresolved\_Regime\_Characterization;
Arc\_QC\_3\_Failure\_Mode\_Rates;
Arc\_QC\_4\_Architectural\_Class\_Audit;
Arc\_QC\_5\_Multiplicity\_Cap\_Function;
Arc\_QC\_6\_Predictive\_Content; Arc\_QC\_7\_Synthesis.

\textbf{Related ED Papers.} \emph{Black Holes in Event Density: A
Substrate-Level Architecture} (May 2026); \emph{Foundations of Substrate
Gravity: From Newtonian Limit to Curvature Emergence} (April 2026);
\emph{Event Density Foundations: A Unified Substrate Architecture for
Quantum, Fluid, Gauge, and Gravitational Dynamics} (April 2026);
\emph{Navier-Stokes Synthesis Paper} (April 2026, with Appendices C/D/E
covering MHD, DCGT, Yang-Mills); \emph{Theorem 17:
Gauge-Field-as-Rule-Type}; \emph{Theorem 18: V1 Kernel Retardation}.

\textbf{ED-I Series Foundations.} ED-I-01 (Superconductivity); ED-I-12
(Photonics: Yablonovitch, Pendry, Capasso); ED-I-13 (Quantum
Information: Bennett, Brassard, Deutsch, Jozsa, Shor); ED-I-14
(Topological Effects: Aharonov, Berry, Bohm, Casher, Tonomura); ED-I-18
(Multi-Timescale Photonics: Hafezi 2025); ED-I-23 (Josephson Junction
Physics); ED-I-29 (Tunneling).

\textbf{Citation Block (Zenodo-safe).}

\begin{quote}
Proxmire, A. (2026). \emph{Quantum Computation in Event Density: A
Substrate-Level Architecture and Its Limits.} Event Density Program,
Quantum Computing Foundations Paper, May 2026.
\end{quote}

\textbf{Conditional-Positive Verdict Disclaimer.} This paper presents a
substrate-level architectural account of quantum computation.
Form-derivations are FORCED from substrate machinery; value-layer
details (\(\mathcal{M}_{\mathrm{crit}}\), \(N_\star\),
\(N_{\mathrm{corr}}\), \(\Delta_{\mathrm{top}}^{\max}\), specific shapes
of \(\mu, \kappa, g/h/c\)) are INHERITED and not closed-form derived at
present. No constructive recipe for fault-tolerant QC is delivered; the
arc is structural, not engineering. Open extensions (closed-form
\(\mathcal{M}_{\mathrm{crit}}\), architecture-to-platform calibration,
topology-stability theorem, surface-code recursive-overlay derivation,
hybrid-handoff dynamics) are explicitly named.

\end{document}
